\documentclass[11pt, a4paper]{article}
\usepackage[utf8]{inputenc}
\usepackage{geometry}
\geometry{margin=1in}
\usepackage{graphicx}
\usepackage{amsmath}
\usepackage{booktabs}
\usepackage{longtable}
\usepackage{xcolor}
\usepackage{enumitem}
\usepackage{fancyhdr}
\usepackage{titlesec}
\usepackage{hyperref}
\hypersetup{colorlinks=true, linkcolor=blue!70!black, urlcolor=blue!70!black}

% Header/Footer
\pagestyle{fancy}
\fancyhf{}
\fancyhead[L]{\textsc{TWP Logbook — Assignment 03}}
\fancyhead[R]{\textsc{CSP529}}
\fancyfoot[C]{\thepage}
\renewcommand{\headrulewidth}{0.4pt}

\titleformat{\section}{\large\bfseries}{\thesection.}{0.5em}{}[\titlerule]
\titleformat{\subsection}{\normalsize\bfseries}{\thesubsection.}{0.5em}{}

\title{
    \vspace{-1cm}
    \textbf{Writing Logbook --- Assignment 03} \\[0.3em]
    \large CSP529 Technical Writing \& Publishing \\[0.2em]
    \large \textit{From Reading to Evidence: What Am I Really Saying?}
}
\author{
    \textbf{MT26MCS024} --- Gaurav Acharya \quad
    \textbf{MT26MCS032} --- Sandeep
}
\date{September 2026}

\begin{document}
\maketitle
\thispagestyle{fancy}

% ============================================================
\section{Assignment Overview}
% ============================================================

\begin{tabular}{@{}p{4.5cm} p{10cm}@{}}
    \toprule
    \textbf{Item} & \textbf{Detail} \\
    \midrule
    Practical Title & Practical 3 --- From Reading to Evidence: What Am I Really Saying? \\
    Article Analysed & ``System Correctness Practices at Amazon Web Services'' \\
    Group Members & Gaurav Acharya (MT26MCS024), Sandeep (MT26MCS032) \\
    Deliverables & Article Summarization (400 words), Classification Table, Original Article Analysis \\
    Date of Completion & September 2026 \\
    \bottomrule
\end{tabular}

% ============================================================
\section{Objectives \& Learning Outcomes}
% ============================================================

This assignment focused on analytical reading and evidence-based reasoning in technical writing. Before learning \emph{how} to cite sources, we first learned to ask: \emph{What exactly am I saying, and why should the reader believe it?} The learning outcomes included:

\begin{itemize}[nosep]
    \item Reading a professional technical article analytically rather than passively.
    \item Identifying the principal message and supporting ideas.
    \item Summarizing technical material accurately in one's own words (max 400 words).
    \item Classifying every sentence by type: Fact (F), Definition (D), Claim (C), Observation (O), Inference/Interpretation (I), Assumption/Prediction (A).
    \item Determining which statements require supporting evidence and which can stand without citation.
    \item Comparing professional writing patterns with one's own writing.
\end{itemize}

% ============================================================
\section{Work Performed --- Stage by Stage}
% ============================================================

\subsection{Stage 1: Read Before You Write}
We read the complete article on System Correctness Practices at AWS without using any AI tools. During reading, we identified:
\begin{itemize}[nosep]
    \item \textbf{Problem/topic:} Maintaining system correctness in large-scale distributed cloud systems.
    \item \textbf{Importance:} Bugs in cloud infrastructure can affect millions of users; formal methods can prevent catastrophic failures.
    \item \textbf{Principal message:} AWS combines heavyweight formal methods (TLA+, P Language) with lightweight approaches (property-based testing, deterministic simulation, continuous fuzzing) to ensure system correctness.
    \item \textbf{Evidence presented:} Migration of S3 to read-after-write consistency validated via P Language; PObserve tool for production validation; 733 FIS fault-injection experiments for Prime Day 2024.
    \item \textbf{Conclusions:} Formal methods are crucial for cloud system safety but face adoption barriers (steep learning curve, specialized expertise, poor UI).
    \item \textbf{Uncertainties:} Whether increased industry adoption will attract sufficient talent; metastability failure modelling remains expensive.
\end{itemize}

\subsection{Stage 2: Summarize (400 Words)}
We produced a structured summary titled \textit{``System Correctness Practices at Amazon Web Services''} covering:
\begin{enumerate}[nosep]
    \item TLA+ as the historical correctness tool and P Language as a more accessible alternative.
    \item Validation achievements: S3 migration, PObserve tool (2023).
    \item Lightweight formal methods: property-based testing, deterministic simulation, continuous fuzzing.
    \item AWS Fault Injection Service and metastable failure challenges.
    \item Performance optimization via HOL Light theorem prover.
    \item Adoption barriers and educational challenges.
\end{enumerate}

\subsection{Stage 3: Sentence Classification Table}
We classified every sentence in our summary. Key classifications included:

\begin{longtable}{@{}c p{6cm} c c p{4cm}@{}}
    \toprule
    \textbf{No.} & \textbf{Sentence (abbreviated)} & \textbf{Type} & \textbf{Evidence?} & \textbf{Reasoning} \\
    \midrule
    \endhead
    1 & TLA+ was used historically by AWS for system correctness\ldots & F & Yes & Verifiable historical fact, but requires source \\
    2 & P Language sought a balance between math capabilities and ease of understanding\ldots & C & Yes & ``Easier to understand'' is subjective \\
    3 & S3 migrated to read-after-write consistency using P\ldots & C & Yes & Claim requiring proof of actual tool usage \\
    4 & PObserve was built in 2023\ldots & F & Yes & Verifiable date/tool release \\
    5 & Property-based testing combines\ldots & D & No & Defines terminology \\
    6 & Deterministic simulation executes on single thread\ldots & D & No & Defines a concept \\
    7 & Continuous fuzzing covers scalable test-input generation\ldots & D & No & Defines a concept \\
    8 & AWS launched Fault Injection Service\ldots & F & Yes & Service launch needs verification \\
    9 & Formal methods are a crucial part of safely improving\ldots & I & Yes & ``Crucial'' is a subjective judgment \\
    10 & Primary barriers are steep learning curve\ldots & C & Yes & Assertion requiring empirical support \\
    \bottomrule
\end{longtable}

\noindent\textbf{Key Insight:} Definitions (D) typically do not require citations because they establish terminology. Facts (F) and Claims (C) require evidence. Inferences (I) represent the author's conclusions and need the strongest evidentiary support.

\subsection{Stage 4: The Evidence Question}
For every sentence, we asked \textit{``How do I know this?''} and determined whether evidence was required. We discovered that:
\begin{itemize}[nosep]
    \item Not every sentence needs a citation, but every substantive statement must be defensible.
    \item Definitions can stand without citation in context, but facts about specific corporate practices always need a source.
\end{itemize}

\subsection{Stage 5: Difficult Sentences}
We identified two sentences that were most difficult to classify:

\begin{enumerate}[nosep]
    \item \textbf{``P Language sought a balance\ldots''} --- Classified as \textbf{Claim (C)} rather than Fact (F) because ``easier to understand'' is inherently subjective and requires empirical support. Could alternatively be Inference (I) since the authors are interpreting the designers' intent.
    
    \item \textbf{``Formal methods are a crucial part\ldots''} --- Classified as \textbf{Inference (I)} rather than Fact (F) because ``crucial'' represents a value judgment derived from the evidence rather than an independently verifiable truth.
\end{enumerate}

\subsection{Stage 6: Examining the Professional Author}
We selected six sentences from the original article and classified them:

\begin{longtable}{@{}p{5cm} c c c p{4.5cm}@{}}
    \toprule
    \textbf{Original Sentence} & \textbf{Type} & \textbf{Ev.\ Req?} & \textbf{Ev.\ Given?} & \textbf{Reasoning} \\
    \midrule
    \endhead
    Modelling the Aurora commit protocol in P and TLA+\ldots & I & Yes & Yes & Analytical deduction from model checking \\
    Amazon ran 733 FIS experiments for Prime Day 2024 & O & Yes & Yes & Concrete empirical metric \\
    We anticipate increased industry adoption\ldots & A & Yes & No & Forward-looking speculation (``We anticipate'') \\
    \bottomrule
\end{longtable}

\noindent\textbf{Key difference observed:} The professional author consistently signals the strength of each statement (``we anticipate,'' ``we were able to prove,'' ``this allowed us to identify''), whereas our summary tended to present inferences as if they were established facts.

\subsection{Stage 7: Strength of a Claim}
We analysed the strongest claim in our summary and explored how evidence quality should calibrate claim strength. This directly prepared us for the claim-evidence tracing work in Assignment~04.

% ============================================================
\section{Challenges Encountered}
% ============================================================

\begin{enumerate}[nosep]
    \item \textbf{Classification ambiguity:} Many sentences could reasonably belong to multiple categories (e.g., Fact vs.\ Claim, Claim vs.\ Inference). The exercise forced us to articulate \emph{why} one classification was more appropriate.
    \item \textbf{Summarization discipline:} Condensing a multi-page technical article into 400 words required difficult decisions about what to omit.
    \item \textbf{Recognizing hedging language:} We initially missed professional hedging signals (``we anticipate,'' ``may sometimes'') that indicate assumptions vs.\ facts.
    \item \textbf{\LaTeX\ table formatting:} The classification table with wrapped text in columns required careful use of \texttt{p\{\}} column types and \texttt{longtable} for multi-page tables.
\end{enumerate}

% ============================================================
\section{Key Takeaways}
% ============================================================

\begin{enumerate}[nosep]
    \item The question ``What kind of statement am I making?'' must come \emph{before} the question ``How do I cite this?''
    \item Not every sentence needs a citation, but every substantive statement must be defensible.
    \item Professional writers carefully signal statement strength through hedging language --- a skill student writers must develop.
    \item The distinction between ``sounds like a fact'' and ``is actually a claim'' is subtle but critically important for academic integrity.
    \item Reading analytically (looking for statement types, evidence quality, and hedging) is fundamentally different from reading passively (just absorbing information).
\end{enumerate}

\end{document}
